\documentclass[aspectratio=169,10pt]{beamer}

\usepackage[T1]{fontenc}
\usepackage{lmodern}

\definecolor{Navy}{HTML}{172D45}
\definecolor{Teal}{HTML}{087F8C}
\definecolor{Muted}{HTML}{506175}
\definecolor{Amber}{HTML}{A65420}
\definecolor{RuleGrey}{HTML}{D8E0E7}

\setbeamercolor{normal text}{fg=Navy,bg=white}
\setbeamercolor{structure}{fg=Teal}
\setbeamercolor{frametitle}{fg=Navy,bg=white}
\setbeamercolor{alerted text}{fg=Amber}
\setbeamertemplate{navigation symbols}{}
\setbeamertemplate{frametitle}{%
  \vspace{2mm}\insertframetitle\par\vspace{1mm}%
  {\color{Teal}\rule{\textwidth}{0.65pt}}\vspace{-1mm}}
\setbeamersize{text margin left=9mm,text margin right=9mm}
\setbeamertemplate{footline}{%
  \hfill\scriptsize
  \insertframenumber\,/\,\inserttotalframenumber\hspace{9mm}\vspace{2mm}}

\newcommand{\source}[1]{\vfill{\fontsize{6.2}{7.5}\selectfont\color{Muted}Source: #1\par}}
\newcommand{\seqitem}[3]{%
  {\normalsize\bfseries\color{Teal}#1\quad\textcolor{Navy}{#2}}\par
  \vspace{0.6mm}
  {\small #3}\par}
\newcommand{\seqrule}{%
  \vspace{1mm}{\color{RuleGrey}\rule{\textwidth}{0.45pt}}\vspace{1mm}}

\begin{document}

\begin{frame}[t]{The decision after acceptance rejection}
  \vspace{1mm}
  {\Large\bfseries Every stack in scope has failed required acceptance testing.}\par
  \vspace{4mm}

  \seqitem{01}{Acceptance result}
    {The HBM3E stack is rejected. The diagnostic path does not change its disposition.}

  \seqrule

  \seqitem{02}{Diagnostic decision}
    {The on-duty quality engineer chooses which of five diagnostic procedures to run next.}

  \seqrule

  \seqitem{03}{Process value}
    {The findings identify fault mechanisms for process engineering. The stack will be scrapped either way.}

  \vspace{3mm}
  {\small\bfseries\color{Amber}Helion supports the investigation. The engineer owns the decision.}

  \source{Client brief, Business Context and Objective.}
\end{frame}

\begin{frame}[t]{The investigation before a procedure}
  \vspace{2mm}

  \seqitem{01}{Load available records}
    {Use manufacturing, package-inspection and acceptance observations available at the decision time.}

  \seqrule

  \seqitem{02}{Establish required work}
    {Open diagnostic concerns, then apply procedure eligibility and evidence-preservation rules.}

  \seqrule

  \seqitem{03}{Choose the next procedure}
    {The candidate baseline uses a fixed rule priority. Helion ranks the same eligible procedures using seven initial fault probabilities.}

  \vspace{3mm}
  {\small\bfseries\color{Amber}The ML proposal changes the ranking step only.}

  \source{MOCK-ENG-002, decision boundary and evidence record. Design notebook, policy section.}
\end{frame}

\begin{frame}[t]{Evidence update and closure}
  \vspace{2mm}
  {\Large\bfseries Each procedure changes the investigation state.}\par
  \vspace{4mm}

  \seqitem{01}{Run the procedure}
    {The engineer approves the recommendation before the lab begins work.}

  \seqrule

  \seqitem{02}{Record the finding}
    {Record each scoped result as confirmed present, confirmed absent or inconclusive.}

  \seqrule

  \seqitem{03}{Update and repeat}
    {Refresh unresolved obligations and procedure eligibility, then recommend the next action.}

  \seqrule

  \seqitem{04}{Approve closure}
    {The engineer closes the case only after the required diagnostic concerns have been resolved.}

  \source{MOCK-ENG-002, evidence record and stopping logic.}
\end{frame}

\begin{frame}[t]{Business objective and diagnostic guardrail}
  \vspace{1mm}
  {\Large\bfseries Lower cost only counts at the same diagnostic standard.}\par
  \vspace{1mm}

  \seqitem{01}{Objective}
    {Reduce the standard resources consumed per started investigation relative to the rule baseline.}

  \seqrule

  \seqitem{02}{Cost scope}
    {Count engineer and technician time, equipment occupancy and procedure supplies.}

  \seqrule

  \seqitem{03}{Quality guardrail}
    {Match the rule baseline on diagnostic completeness and false-negative performance.}

  \seqrule

  \seqitem{04}{Reporting meaning}
    {Track turnaround separately. Resource dollars measure consumed capacity, not demonstrated avoidable cash.}


  \source{Client brief, Constraints and scope. SYN-OPS-001 cost-accounting assumptions.}
\end{frame}

\begin{frame}[t]{From observations to one recommendation}
  \vspace{2mm}

  \seqitem{01}{Inputs}
    {Use initial manufacturing, package-inspection and acceptance observations.}

  \seqrule

  \seqitem{02}{Model output}
    {Produce seven independent fault probabilities so that mechanisms may coexist.}

  \seqrule

  \seqitem{03}{Decision policy}
    {Combine probabilities with procedure coverage, inconclusive rates and standard resource costs.}

  \seqrule

  \seqitem{04}{Recommendation}
    {Rank only the currently eligible procedures and present the next one for engineer review.}

  \vspace{3mm}
  {\small\bfseries\color{Amber}One model supports one next-test decision.}

  \source{Client brief, The Ask and Constraints and scope. Design notebook, policy definition.}
\end{frame}

\begin{frame}[t]{Safety boundary and human control}
  \vspace{2mm}

  \seqitem{01}{Engineer authority}
    {The engineer approves each procedure and the eventual closure decision.}

  \seqrule

  \seqitem{02}{No automatic absence}
    {Low probability, an unexamined mechanism and an inconclusive result never confirm absence.}

  \seqrule

  \seqitem{03}{Coexisting faults}
    {A first confirmed mechanism does not cancel other open diagnostic concerns.}

  \seqrule

  \seqitem{04}{Evidence updates}
    {New findings change mechanism states and procedure eligibility. Initial model probabilities remain fixed.}

  \vspace{3mm}
  {\small\bfseries\color{Amber}The baseline and ML route retain the same diagnostic obligations.}

  \source{MOCK-ENG-002. Design notebook, policy and operating sections.}
\end{frame}

\end{document}
